% Phần đóng góp: Nguyễn Thành Bảo. Đối chiếu BackEnd/Solver/UCSSolver.py.
\section{Tìm kiếm theo chi phí đồng nhất (UCS)}
\label{sec:ucs}

\subsection{Mô hình trạng thái và chi phí}

Thuật toán UCS (Uniform-Cost Search) tìm lời giải có tổng chi phí nhỏ nhất.
Trong chương trình, mỗi đỉnh của đồ thị là một \texttt{GameState}, gồm trạng
thái khối và tập cầu đang mở. Trạng thái khối thể hiện vị trí, hướng đặt,
trạng thái tách đôi và khối con đang được điều khiển. Vì vậy, hai trạng thái
cùng vị trí nhưng khác cấu hình cầu hoặc khối đang hoạt động vẫn có thể có
các bước đi hợp lệ khác nhau.

Các cạnh được sinh qua \texttt{GameRules.apply\_action()}, dùng chung luật chơi
của nhóm. Tập hành động gồm lên, xuống, trái, phải và đổi khối con
(\texttt{SWITCH}). Chỉ những hành động tạo được trạng thái hợp lệ mới trở
thành cạnh. Trong phiên bản này, mọi hành động có chi phí bằng 1, kể cả đổi
khối con. Với đường đi gồm $k$ hành động:
\[
    g(n)=\sum_{i=1}^{k}c_i=k.
\]
Do đó, tối thiểu hóa chi phí cũng là tối thiểu hóa số hành động.

\subsection{Cấu trúc dữ liệu và quá trình tìm kiếm}

Lớp \texttt{UCSSolver} kế thừa \texttt{BaseSolver}, dùng min-heap của Python
để lấy node có chi phí nhỏ nhất. Mỗi phần tử heap có dạng
$(g,\,\text{thứ tự thêm},\,\text{node})$. Bộ đếm thứ tự giúp xử lý trường hợp
đồng chi phí mà không phải so sánh trực tiếp các đối tượng trạng thái.

Bảng \texttt{best\_cost} ánh xạ toàn bộ trạng thái tới chi phí thấp nhất đã
biết. Khi tìm được đường đi tốt hơn, thuật toán cập nhật bảng và thêm một
phần tử mới vào heap. Phần tử cũ vẫn có thể còn trong heap; khi lấy ra, nó
được bỏ qua nếu chi phí không còn bằng giá trị tốt nhất. Thuật toán chỉ
kiểm tra đích sau bước loại phần tử cũ này, khi node được lấy khỏi heap.
Lời giải được truy vết qua các trường \texttt{parent} và \texttt{action}
của \texttt{SearchNode}~\cite{bloxorz_ucs_source}.

\begin{algorithm}[htbp]
    \caption{UCS trên trạng thái Bloxorz}
    \label{alg:ucs}
    \begin{algorithmic}[1]
        \Require Trạng thái đầu $s_0$, bộ luật sinh trạng thái kế tiếp
        \Ensure Chuỗi hành động tối ưu hoặc không có lời giải
        \State $Q\gets$ min-heap chứa node gốc với chi phí $0$
        \State $best[s_0]\gets 0$; các giá trị khác mặc định $+\infty$
        \While{$Q$ chưa rỗng}
            \State Lấy node $n$ có chi phí nhỏ nhất khỏi $Q$
            \If{$g(n)\ne best[state(n)]$}
                \State Bỏ qua $n$ và tiếp tục vòng lặp
            \Else
                \If{$state(n)$ là đích}
                    \State \Return chuỗi hành động truy vết từ $n$
                \EndIf
                \For{mỗi $(a,s',c)$ hợp lệ từ $state(n)$}
                    \State $g'\gets g(n)+c$
                    \If{$g'<best[s']$}
                        \State $best[s']\gets g'$
                        \State Thêm node con $(s',n,a,g')$ vào $Q$
                    \EndIf
                \EndFor
            \EndIf
        \EndWhile
        \State \Return không có lời giải
    \end{algorithmic}
\end{algorithm}

\subsection{Tính chất và giới hạn}

Với số trạng thái hữu hạn và chi phí cạnh dương, khi lấy node hợp lệ có
chi phí nhỏ nhất ra khỏi heap, không thể có đường đi chưa xử lý nào tới
node đó với chi phí thấp hơn: đường đi như vậy phải đi qua một node có
chi phí nhỏ hơn, vốn đã được ưu tiên xử lý trước. Vì vậy, khi node đích
được lấy ra, chi phí lời giải là tối ưu. Nếu heap rỗng, toàn bộ phần đồ thị
có thể đi tới đã được duyệt mà không tìm thấy đích.

Gọi $V$ và $E$ là số trạng thái và cạnh trong đồ thị có thể đi tới. Nếu xem
chi phí sinh trạng thái và tra bảng băm là hằng số, cách cài đặt heap cho
phép giữ phần tử cũ có thời gian $O((V+E)\log(V+E))$ và bộ nhớ phụ
$O(V+E)$ trong trường hợp xấu nhất. Các giới hạn thực tế của chương trình
có thể dừng tìm kiếm trước khi thuật toán kết luận; trường hợp này được
báo là vượt giới hạn, không phải vô nghiệm.

Vì mọi hành động hiện có cùng chi phí, BFS và UCS đều tối thiểu hóa số
bước. UCS không được kỳ vọng luôn nhanh hơn BFS; heap còn có chi phí quản
lý riêng. Lợi ích của cơ chế cập nhật chi phí thể hiện khi mở rộng bài toán
sang các cạnh có trọng số khác nhau. Trường hợp này đã được kiểm thử bằng
đồ thị tổng hợp, không thay đổi hàm chi phí của game.
